\section{Использование свойств ядра}
Для каждого наблюдения \( \Xv_i \) и центра \( \xv_j \) определим
вектор разности
\begin{EQA}[c]
     D_{ij}
    =
    \Xv_i - \xv_j
    \in \R^{\dimd}.
\end{EQA}
Соответствующий локализующий вес задаётся формулой
\begin{EQA}[c]
    \weight_{ij}
    =
    \Kern\bigl(
        \| \TEDR  D_{ij} \|^2
    \bigr).
\end{EQA}

В single-index случае локализующий тензор определяется соотношением
\begin{EQA}[c]
    \TEDR^2
    =
    \hiso^{-2}
    \bigl(
        \aniso^2 \Id_{\dimd}
        +
        \betav \betav^{\T}
    \bigr).
\end{EQA}
Поэтому соответствующую квадратичную форму можно вычислить напрямую:
\begin{EQA}[rcl]
    \| \TEDR  D_{ij} \|^2
    &=&
     D_{ij}^{\T}
    \TEDR^2
     D_{ij}
    \\[1mm]
    &=&
    \hiso^{-2}
    \left(
        \aniso^2 \|  D_{ij} \|^2
        +
        \bigl\langle \betav, D_{ij} \bigr\rangle^2
    \right).
\end{EQA}
Следовательно, вес можно записать в векторной форме без явного
построения матрицы \( \TEDR \):
\begin{EQA}[c]
    \weight_{ij}
    =
    \Kern\left(
        \frac{
            \aniso^2 \|  D_{ij} \|^2
            +
            \bigl\langle \betav, D_{ij} \bigr\rangle^2
        }{
            \hiso^2
        }
    \right).
\end{EQA}
Для используемого ядра
\( \Kern(t)=(1-t^2)_{+} \) получаем
\begin{EQA}[c]
    \weight_{ij}
    =
    \left[
        1
        -
        \left(
            \frac{
                \aniso^2 \|  D_{ij} \|^2
                +
                \bigl\langle \betav, D_{ij} \bigr\rangle^2
            }{
                \hiso^2
            }
        \right)^2
    \right]_{+}.
\end{EQA}
В multi-index случае структурная матрица имеет вид
\begin{EQA}[c]
    \JEDR_{k-1}
    =
    \PEDR_{k-1}^{\T}
    \Lambda_{k-1}
    \PEDR_{k-1},
    \qquad
    \PEDR_{k-1}\PEDR_{k-1}^{\T}
    =
    \Id_{\dimm},
\end{EQA}
где
\(
\PEDR_{k-1}\in\R^{\dimm\times\dimd}
\)
задаёт оценённое EDR-подпространство, а
\(
\Lambda_{k-1}\in\R^{\dimm\times\dimm}
\)
содержит веса его направлений.

Для основного варианта локализующего тензора
\begin{EQA}[c]
    \TEDR_k^2
    =
    \aniso_k^2 \hiso_k^{-2}
    \bigl(
        \Id_{\dimd}
        -
        \PEDR_{k-1}^{\T}\PEDR_{k-1}
    \bigr)
    +
    \hiso_k^{-2}\JEDR_{k-1}
\end{EQA}
его квадратичная форма может быть вычислена напрямую:
\begin{EQA}[rcl]
    \| \TEDR_k  D_{ij} \|^2
    &=&
     D_{ij}^{\T}
    \TEDR_k^2
     D_{ij}
    \\[1mm]
    &=&
    \hiso_k^{-2}
    \Bigl[
        \aniso_k^2
        \bigl(
            \| D_{ij}\|^2
            -
            \|\PEDR_{k-1} D_{ij}\|^2
        \bigr)
        +
        \|
            \Lambda_{k-1}^{1/2}
            \PEDR_{k-1} D_{ij}
        \|^2
    \Bigr].
\end{EQA}

Следовательно, вес можно вычислять в векторной форме без явного
построения матрицы \( \TEDR_k\in\R^{\dimd\times\dimd} \):
\begin{EQA}[c]
    \boxed{
    \weight_{ij}^{(k)}
    =
    \Kern\left(
        \frac{
            \aniso_k^2
            \left(
                \| D_{ij}\|^2
                -
                \|\PEDR_{k-1} D_{ij}\|^2
            \right)
            +
            \|
                \Lambda_{k-1}^{1/2}
                \PEDR_{k-1} D_{ij}
            \|^2
        }{
            \hiso_k^2
        }
    \right).
    }
\end{EQA}

Для используемого ядра
\( \Kern(t)=(1-t^2)_+ \) получаем
\begin{EQA}[c]
    \weight_{ij}^{(k)}
    =
    \left[
        1 -
        \left(
            \frac{
                \aniso_k^2
                \left(
                    \| D_{ij}\|^2
                    -
                    \|\PEDR_{k-1} D_{ij}\|^2
                \right)
                +
                \|
                    \Lambda_{k-1}^{1/2}
                    \PEDR_{k-1} D_{ij}
                \|^2
            }{
                \hiso_k^2
            }
        \right)^2
    \right]_+ .
\end{EQA}

\subsection{Критерий нуль веса}

Для используемого компактно поддержанного ядра
\[
\Kern(t)=(1-t^2)_+,
\]
где аргумент \(t\geq 0\), имеем
\begin{EQA}[c]
    \Kern(t)=0
    \quad\Longleftrightarrow\quad
    t\geq 1.
\end{EQA}
Следовательно, для веса
\begin{EQA}[c]
    \weight_{ij}
    =
    \Kern\bigl(
        \|\TEDR D_{ij}\|^2
    \bigr)
\end{EQA}
получаем общий критерий
\begin{EQA}[c]
        \weight_{ij}=0
        \quad\Longleftrightarrow\quad
        \|\TEDR D_{ij}\|^2\geq 1.
\end{EQA}

\subsubsection{Single-index случай}

В single-index случае
\begin{EQA}[c]
    \|\TEDR D_{ij}\|^2
    =
    \frac{
        \aniso^2\| D_{ij}\|^2
        +
        \langle\betav, D_{ij}\rangle^2
    }{
        \hiso^2
    }.
\end{EQA}
Поэтому
\begin{EQA}[rcl]
    \weight_{ij}=0
    &\Longleftrightarrow&
    \frac{
        \aniso^2\| D_{ij}\|^2
        +
        \langle\betav, D_{ij}\rangle^2
    }{
        \hiso^2
    }
    \geq 1
    \\[1mm]
    &\Longleftrightarrow&
    \aniso^2\| D_{ij}\|^2
    +
    \langle\betav, D_{ij}\rangle^2
    \geq
    \hiso^2.
\end{EQA}
Таким образом, точный критерий нулевого веса имеет вид
\begin{EQA}[c]
        \aniso^2\| D_{ij}\|^2
        +
        \langle\betav, D_{ij}\rangle^2
        \geq
        \hiso^2
        \quad\Longleftrightarrow\quad
        \weight_{ij}=0.
\end{EQA}

Из него можно получить более дешёвый достаточный критерий, использующий
только проекцию на текущее EDR-направление. Поскольку
\begin{EQA}[c]
    \| D_{ij}\|^2
    \geq
    \langle\betav, D_{ij}\rangle^2,
    \qquad
    \|\betav\|=1,
\end{EQA}
имеем
\begin{EQA}[rcl]
    \aniso^2\| D_{ij}\|^2
    +
    \langle\betav, D_{ij}\rangle^2
    &\geq&
    (1+\aniso^2)
    \langle\betav, D_{ij}\rangle^2.
\end{EQA}
Следовательно,
\begin{EQA}[c]
    (1+\aniso^2)
    \langle\betav, D_{ij}\rangle^2
    \geq
    \hiso^2
\end{EQA}
гарантирует нулевой вес. Эквивалентно,
\begin{EQA}[c]
    \boxed{
        \left|
            \langle\betav, D_{ij}\rangle
        \right|
        \geq
        \frac{\hiso}{\sqrt{1+\aniso^2}}
        \quad\Longrightarrow\quad
        \weight_{ij}=0.
    }
\end{EQA}
Этот критерий позволяет отбрасывать пары \((i,j)\) после вычисления
одной скалярной проекции, не вычисляя полную норму
\(\| D_{ij}\|^2\).

\subsubsection{Multi-index случай}

В multi-index случае
\begin{EQA}[c]
    \|\TEDR_k D_{ij}\|^2
    =
    \frac{
        \aniso_k^2
        \left(
            \| D_{ij}\|^2
            -
            \|\PEDR_{k-1} D_{ij}\|^2
        \right)
        +
        \left\|
            \Lambda_{k-1}^{1/2}
            \PEDR_{k-1} D_{ij}
        \right\|^2
    }{
        \hiso_k^2
    }.
\end{EQA}
Следовательно,
\begin{EQA}[rcl]
    \weight_{ij}^{(k)}=0
    &\Longleftrightarrow&
    \aniso_k^2
    \left(
        \| D_{ij}\|^2
        -
        \|\PEDR_{k-1} D_{ij}\|^2
    \right)
    +
    \left\|
        \Lambda_{k-1}^{1/2}
        \PEDR_{k-1} D_{ij}
    \right\|^2
    \geq
    \hiso_k^2.
\end{EQA}
Таким образом, точный критерий имеет вид
\begin{EQA}[c]
    \begin{aligned}
        &
        \aniso_k^2
        \left(
            \| D_{ij}\|^2
            -
            \|\PEDR_{k-1} D_{ij}\|^2
        \right)
        \\
        &\qquad+
        \left\|
            \Lambda_{k-1}^{1/2}
            \PEDR_{k-1} D_{ij}
        \right\|^2
        \geq
        \hiso_k^2
        \quad\Longleftrightarrow\quad
        \weight_{ij}^{(k)}=0.
    \end{aligned}
\end{EQA}

Поскольку строки \( \PEDR_{k-1} \) ортонормированы,
\(\PEDR_{k-1}^{\T}\PEDR_{k-1}\) представляет ортогональный проектор.
Поэтому
\begin{EQA}[c]
    \| D_{ij}\|^2
    -
    \|\PEDR_{k-1} D_{ij}\|^2
    \geq 0.
\end{EQA}
Отбрасывая этот неотрицательный член, получаем нижнюю оценку
\begin{EQA}[c]
    \|\TEDR_k D_{ij}\|^2
    \geq
    \frac{
        \left\|
            \Lambda_{k-1}^{1/2}
            \PEDR_{k-1} D_{ij}
        \right\|^2
    }{
        \hiso_k^2
    }.
\end{EQA}
Следовательно, достаточно проверить условие
\begin{EQA}[c]
    \boxed{
        \left\|
            \Lambda_{k-1}^{1/2}
            \PEDR_{k-1} D_{ij}
        \right\|^2
        \geq
        \hiso_k^2
        \quad\Longrightarrow\quad
        \weight_{ij}^{(k)}=0.
    }
\end{EQA}
Или, эквивалентно,
\begin{EQA}[c]
        (\PEDR_{k-1} D_{ij})^{\T}
        \Lambda_{k-1}
        (\PEDR_{k-1} D_{ij})
        \geq
        \hiso_k^2
        \quad\Longrightarrow\quad
        \weight_{ij}^{(k)}=0.
\end{EQA}

В отличие от точного критерия, последнее условие является только
достаточным: если оно не выполняется, вес всё ещё может оказаться
нулевым за счёт ортогональной составляющей
\[
\aniso_k^2
\left(
    \| D_{ij}\|^2
    -
    \|\PEDR_{k-1} D_{ij}\|^2
\right).
\]
При этом проверка достаточного критерия выполняется целиком в
\(\dimm\)-мерном EDR-пространстве и не требует вычисления полной
\(\dimd\)-мерной квадратичной формы.


\subsection{Выбор $\rho_k$ при фиксированном уменьшении ширины окна}

Аргумент ядра в single-index и multi-index случаях можно записать в общей форме
\begin{EQA}[c]
    q_{ij}(z,h)
    =
    \frac{A_{ij}z+B_{ij}}{h^2},
    \qquad
    z=\rho^2,
    \qquad
    A_{ij},B_{ij}\geq0.
\end{EQA}

Для single-index случая
\begin{EQA}[c]
    A_{ij}=\| D_{ij}\|^2,
    \qquad
    B_{ij}
    =
    \langle\betav_{k-1}, D_{ij}\rangle^2.
\end{EQA}

Для multi-index случая, если
\[
    \uv_{ij}=\PEDR_{k-1} D_{ij},
\]
то
\begin{EQA}[c]
    A_{ij}
    =
    \| D_{ij}\|^2-\|\uv_{ij}\|^2,
    \qquad
    B_{ij}
    =
    \uv_{ij}^{\T}\Lambda_{k-1}\uv_{ij}.
\end{EQA}

Определим суммарную локальную массу
\begin{EQA}[c]
    \mathcal M(\rho,h)
    =
    \sum_{j\in\JJJ}\sumi
    \Kern\left(
        \frac{A_{ij}\rho^2+B_{ij}}{h^2}
    \right),
\end{EQA}
и потребуем
\begin{EQA}[c]
    \mathcal M(\rho,h)\geq \nJJJ\nmin.
\end{EQA}

Поскольку $A_{ij}\geq0$, при увеличении $\rho$ аргумент ядра
не уменьшается, а локальная масса не возрастает. Поэтому при
фиксированном $\hiso_k$ следует выбирать максимально допустимое
значение
\begin{EQA}[c]
    \boxed{
    \rho_k
    =
    \sup
    \left\{
        \rho\in[0,1]:
        \mathcal M(\rho,\hiso_k)
        \geq
        \nJJJ\nmin
    \right\}.
    }
\end{EQA}
В single-index модели параметр $\rho_k$ используется непосредственно, а
в multi-index модели ему соответствует параметр $\alpha_k$. Такой выбор
даёт наиболее узкие окрестности среди всех значений параметра,
удовлетворяющих условию на массу.

Ширина окна уменьшается с фиксированным коэффициентом $a>1$ из конфигурации:
\begin{EQA}[c]
    \hiso_k
    =
    \frac{\hiso_{k-1}}{a}.
\end{EQA}
Если $\hiso_k<\hiso_{\min}$, алгоритм завершается с причиной
\texttt{h\_min}. Иначе при фиксированном $\hiso_k$ находится максимальное
допустимое значение $\rho_k$ по условию на массу выше. Если даже при
$\rho=0$ выполняется
\(
\mathcal M(0,\hiso_k)<\nJJJ\nmin
\), допустимого значения не существует, и алгоритм завершается с причиной
\texttt{local\_mass\_limit}.

Таким образом, на каждом шаге выполняется фактическая последовательность
\[
    \boxed{
    \hiso_k=\frac{\hiso_{k-1}}{a}
    \;\longrightarrow\;
    \text{проверка }\hiso_{\min}
    \;\longrightarrow\;
    \rho_k\ \text{или соответствующий ему }\alpha_k.
    }
\]
Параметры $h_{\mathrm{mass}}$, $\rho_{\min}$ и $a_{\max}$ в этой
реализации не используются; коэффициент $a$ не адаптируется.
